\chapter{Of Records, Porters, and the Post}
\label{art:records}

\articlenote{In which is set down the keeping of records throughout the silo, the carrying of messages by porters where the radio and wire cannot reach, and the obligation that every transaction of consequence to the silo be preserved in the archive for the memory of the silo's own history.}

\section{The Keeping of Records}

Every record the Pact requires to be kept in the Judicial archive under Article 6 is the responsibility of Judicial to keep whole, accurate, and protected from alteration. These records are the silo's memory, and the silo's memory is the silo's greatest treasure after the lives of its citizens.

\begin{clauses}
\item The Judicial archive keeps every ruling, every sentence, every death-finding, every amendment, every standing order of the Mayor, and every record of crime and punishment that the Pact assigns to its keeping.
\item The archive is kept on Level 1, in the Up Top, in a room secured against fire, flood, and interference, and accessible to Judicial officers only, to the Mayor for review of standing orders, and to any citizen who petitions Judicial to see a record concerning that citizen's own business.
\item A citizen may petition to view any record concerning the citizen or the citizen's household, and Judicial shall grant access within one cycle; Judicial may refuse access to records of ongoing investigation or trial, provided the citizen is told the reason for the refusal.
\item The archive is backed by a secondary roll kept in secure keeping by the Head of Supply under Article 11, that if fire or catastrophe should strike the primary archive, the silo's records are not wholly lost.
\end{clauses}

\section{The Porters and the Carrying of Word}

Where a message must travel between floors faster than a citizen might walk, or where a message must be kept private between the sender and the recipient, a porter carries the word. Porters are assigned by the Department of Supply under Article 11 and travel the stairs under protection from the Sheriff.

\begin{clauses}
\item A citizen of majority may become a porter, and porters are trained by Supply in the keeping of messages, the swift travel of the stairs, and the protection of their charges. A porter's shadow under Article 13 shall be one season or two, as the Head of Supply determines.
\item A porter carrying a message has the right of way on the stairs under Article 7; any citizen must yield the porter's passage, and the Sheriff ensures this right is honored.
\item A message carried by a porter is sealed if the sender wishes it sealed, and a porter may not read a sealed message, nor may a citizen intercept a message from a porter without the porter's consent, save the Sheriff or Judicial may do so in pursuit of an investigation. An unsealed message is assumed to be open to ordinary view.
\item A porter who breaks a sealed message is answerable under Article 17 as for breach of trust; a citizen who intercepts a porter's sealed message is answerable as for theft or tampering, as the case requires.
\end{clauses}

\section{The Post and the Radio Communications}

Messages that may travel by radio or wire, as Information Technology under Article 8 maintains those means, travel by radio relay from one floor to the next, faster than a porter might carry them. Information Technology maintains a Post office on the Up Top, and satellite Post rooms on key floors, where citizens may send messages by radio or receive messages sent by radio to them.

\begin{clauses}
\item Any citizen may send a message by radio through the Post office, provided the message is not forbidden speech under Article 16, and the sender pays a fee in chits to Information Technology for the maintenance of the Post service.
\item A message sent by radio is recorded by Information Technology in a log, for the Judicial archive, noting the sender, the recipient, the date, and the time; Information Technology shall not keep the content of private messages, only the fact of their transmission.
\item A citizen may petition for privacy in message-sending, and Information Technology may encrypt a message in a cipher maintained by Information Technology alone, so that no citizen save the recipient may read it; such encryption is the charge of Information Technology and may not be undertaken by any citizen or Department save Information Technology.
\end{clauses}

\section{The Freight Lift Log}

The freight lift under Article 9 is logged in every use, and the log is kept in the Judicial archive for the record of what goods, ordered by what citizen and carried to what floor, have moved through the silo. This log is not open to citizens, but is audited by Judicial to ensure the lift is not being used contrary to Article 9.

\begin{clauses}
\item Every use of the freight lift shall be logged by the operators of the lift under Mechanical's charge: the goods' description and weight, the citizens ordering the lift, the floor of origin, the floor of destination, the date, and the time.
\item The log is kept in duplicate: one copy in the Judicial archive, one in the Mechanical records under Article 9. The two copies are compared once per season by Judicial and the Head of Mechanical to ensure they match and no unauthorized use has been made of the lift.
\end{clauses}

\section{The Amendment Log}

The roll of amendments to this Pact, ratified by the Assembly under Article 20, is kept as a separate record from the main archive, so that citizens may see how the Pact has changed since its founding. This roll is kept in the Judicial archive and in the backmatter of this Pact itself, under Article 23.

\begin{clauses}
\item Every amendment ratified by the Assembly is entered, in order of ratification, in the Amendment Log, with the date of ratification, the text of the amendment, and a statement by the Judge as to whether Judicial found any conflict between the amendment and existing provisions.
\item A citizen may view the Amendment Log in the Judicial archive or may read it directly in the Pact's own backmatter, and may petition Judicial for an explanation of any amendment's text.
\end{clauses}
